\subsection{标架与主纤维化}

12.4.1. 设 \(j \in \mathbf{J}^k(\mathbf{X}, \mathbf{Y})\)，令 \(x = s(j)\)，\(y = b(j)\)。称 \(j\) 可逆，如果存在 \(j' \in \mathbf{J}_y^k(\mathbf{Y}, \mathbf{X})_x\)，使得 \(j' \circ j = j_x^k(\mathrm{Id}_{\mathbf{X}})\) 且 \(j \circ j' = j_y^k(\mathrm{Id}_{\mathbf{Y}})\)；喷射 \(j'\) 因而唯一确定，并记作 \(j^{-1}\)。若 \(k = 0\)，每个喷射均可逆。若 \(k \geqslant 1\)，以下条件等价：

\begin{enumerate}
\def\labelenumi{\alph{enumi})}
\item
  \(j\) 是可逆的；
\item
  映射 \(\mathrm{T}(j)\colon\mathrm{T}_{x}(\mathrm{X})\rightarrow\mathrm{T}_{y}(\mathrm{Y})\)（参见 12.3.4）是同构；
\item
  存在一个同构 \(g\)，类为 \(\mathbf{C}^r\)，从 \(x\) 的开邻域到 \(y\) 的开邻域，使得 \(j_x^k(g) = j\)。
\end{enumerate}

12.4.2. 设 E 是 Banach 空间。令 \(\mathbf{GL}^{k}(\mathbf{E})\) 为从 E 到 E 的、可逆且源和目标均为 0 的 k 阶喷射集合；这是 \(\mathrm{J}_{0}^{k}(\mathrm{E},\mathrm{E})_{0}\) 的开子集。赋予它喷射的复合律以及由 Banach 空间 \(\mathrm{J}_{0}^{k}(\mathrm{E},\mathrm{E})_{0}=\mathrm{Q}^{k}(\mathrm{E},\mathrm{E})\) 诱导的流形结构后，它是一个类 \(C^{\omega}\) 的群流形。

有 \(\mathbf{GL}^0 (\mathbf{E}) = \{e\}\)。群 \(\mathbf{GL}^1 (\mathbf{E})\) 通过 T 与 E 的自同构群 \(\mathbf{GL}(\mathbf{E})\) 等同。

若 \(k' \leqslant k\)，则映射 \(r^{k, k'}: \mathbf{GL}^k(\mathbf{E}) \to \mathbf{GL}^{k'}(\mathbf{E})\) 是类 \(\mathbf{C}^\omega\) 的同态，且是满的浸没。映射 \(f \mapsto \mathrm{Id}_{\mathbf{E}} + f\) 是从 \(\mathrm{P}_k(\mathbf{E}, \mathbf{E})\) 到 \(r^{k, k-1}\) 的核的群流形同构。

12.4.3. 设 E 是 Banach 空间。对于每个 \(x \in X\)，称 X 的 \(k\) 阶 E-标架在 \(x\) 处为 \(J_0^k(E, X)_x\) 中的每个可逆元素。X 的 \(k\) 阶 E-标架集合是开子流形 \(R^k(E, X)\)，位于 \(J_0^k(E, X)\) 中；映射 \(b\) 的限制仍记为 \(b\)，是从 \(R^k(E, X)\) 到 X 的态射；同样，若 \(k' \leqslant k\)，则仍用 \(r^{k, k'}\) 表示从 \(R^k(E, X)\) 到 \(R^{k'}(E, X)\) 的态射，它是 12.1.5 中映射 \(r^{k, k'}\) 的限制。

12.4.4. 假设 X 是纯类型 E（5.1.7）。群 \(\mathbf{GL}^{k}(\mathbf{E})\) 右作用于 \(\mathrm{R}^{k}(\mathrm{E},\mathrm{X})\)，作用律为 \((\rho,u)\mapsto\rho\circ u\)，且四元组 \(\lambda_{\mathbf{X}}=(\mathrm{R}^{k}(\mathbf{E},\mathrm{X}),\mathbf{GL}^{k}(\mathbf{E}),\mathbf{X},\mathbf{b})\) 是主纤维化（6.2.1），结构群为 \(\mathbf{GL}^{k}(\mathbf{E})\)，底空间为 X，类为 \(C^{r-k}\)。

若 \(k' \leqslant k\)，令 H 为 \(r^{k, k'} \colon \mathbf{GL}^{k}(\mathrm{E}) \to \mathbf{GL}^{k'}(\mathrm{E})\) 的核；四元组 \((\mathbf{R}^k(\mathrm{E}, \mathrm{X}), \mathbf{H}, \mathbf{R}^{k'}(\mathrm{E}, \mathrm{X}), r^{k, k'})\) 是类为 \(\mathbf{C}^{r-k}\) 的主纤维化。

流形 \(J^{k}(X,Y)\) 配备了与 \(\lambda_{X}\) 相关联的纤维丛结构（6.5.1）：典型纤维为 \(J_{0}^{k}(E,Y)\)，其上 \(GL^{k}(E)\) 按作用律 \((u,j)\mapsto j\circ u^{-1}\) 左作用；标架映射 \(R^{k}(E,X)\times J_{0}^{k}(E,Y)\to J^{k}(X,Y)\) 将 \((\rho,j)\) 映到 \(j\circ\rho^{-1}\)。与此纤维丛结构相应的投影 \(J^{k}(X,Y)\to X\) 是 s。

12.4.5. 设 F 是 Banach 空间，且 Y 是 F 类型的纯流形。则 \(\mathrm{J}^{k}(\mathrm{X},\mathrm{Y})\) 带有与 \(\lambda_{Y}\) 相关的纤维丛结构：典型纤维为 \(\mathrm{J}^{k}(\mathrm{X},\mathrm{F})_{0}\)，\(\mathrm{GL}^{k}(\mathrm{F})\) 按规律 \((v,j)\mapsto v\circ j\) 左作用于其上；标架映射为 \((\sigma,j)\mapsto\sigma\circ j\)；投影 \(\mathrm{J}^{k}(\mathrm{X},\mathrm{Y})\to\mathrm{Y}\) 是 b。

12.4.6. 取 12.4.4 和 12.4.5 的假设，令 \(\mu\) 为主纤维化

\[
(\mathbf {R} ^ {k} (\mathrm{E}, \mathrm{X}) \times \mathbf {R} ^ {k} (\mathrm{F}, \mathrm{Y}), \mathbf {G L} ^ {k} (\mathrm{E}) \times \mathbf {G L} ^ {k} (\mathrm{F}), \mathrm{X} \times \mathrm{Y}, \boldsymbol {b} \times \boldsymbol {b}),
\]

是 \(\lambda_{\mathbf{X}}\) 与 \(\lambda_{\mathbf{Y}}\) 的乘积。于是，\(\mathrm{J}^{k}(\mathrm{X},\mathrm{Y})\) 配备了与 \(\mu\) 相关联的纤维丛结构：典型纤维为 \(\mathrm{J}_{0}^{k}(\mathrm{E},\mathrm{F})_{0}\)，其上 \(\mathbf{GL}^{k}(\mathrm{E})\times \mathbf{GL}^{k}(\mathrm{F})\) 按作用律 \(((u,v),j)\mapsto v\circ j\circ u^{-1}\) 左作用；标架映射为 \(((\rho ,\sigma),j)\mapsto \sigma \circ j\circ \rho^{-1}\)；投影 \(\mathrm{J}^k (\mathrm{X},\mathrm{Y})\to \mathrm{X}\times \mathrm{Y}\) 为 \((s,b)\)。

